\begin{statement}{C.2.RAT.9}{Proved}{Every rational double ordinary-triple cover has nonnegative Euler number}
Retain the torus-fibre, rational-normalization and specified
ruled-boundary morphism of C.2.RAT.1. Suppose $C$ is not
contracted, its image has generic ordinary-triple transversal,
and $C\to H^\nu$ has degree two. Then $e(S)\ge0$,
including all common zeros and their multiplicities in
the cubic normal tensor.
\end{statement}
\begin{proof}
Again A.5.19 makes $H$ smooth rational. Let $A$ be the
common zero divisor of the cubic coefficients and put
$v=\deg A$. The cases $v=0,1$ are C.2.RAT.7--8.
We first verify that $v\le2$. Write
$N=\OO(a)\oplus\OO(-2-a)$ with $a\ge-1$.
For $a\ge0$ and $v>0$, the coefficient of the cube
of the first normal coordinate has degree $-3a-v$
and vanishes. Thus the other coordinate is a factor.
The remaining quadratic is irreducible over the
function field, because it describes the connected
degree-two cover. Its first-coordinate-square
coefficient must be nonzero and has degree
$2-a-v$. This gives $a+v\le2$.
For $a=-1$, on $\PP^1\times\PP^1$ the section
$R_1$ and quadratic component have bidegrees
$(1,d)$ and $(2,3-v-d)$, with $d\ge0$ and the
first entry denoting degree over $H$. An integral
quadratic component requires $3-v-d\ge1$:
a divisor of bidegree $(2,0)$ is a union of
constant sections over $\C$. Thus $v\le2$ here
as well.

It remains to consider $v=2$. With the notation
of C.2.RAT.7, the horizontal divisor is
\[
 B_h=R_1+R_2\sim3\xi-2F,
       \qquad p_a(B_h)=0,
       \qquad \sum_p\delta_p(B_h)=1.
\tag{C.RD.28}
\]
Both normalization components are rational.
Consequently its only possible singularity is
one ordinary node or one ordinary cusp. An
intersection between $R_1$ and $R_2$, if present,
is the node. Otherwise a node or cusp belongs
to $R_2$. Every other point of the horizontal
divisor is smooth.

\paragraph{Two distinct common zero values.}
First let $A=p+q$ with $p\ne q$. The full branch
divisor is reduced:
$B=R_1+R_2+F_p+F_q\sim3\xi$.
Use the ordinary cyclic algebra of (C.RD.13).
The local analysis in the proof of C.2.RAT.8
applies at each of the two fibres. A nonzero
local geometric genus there requires positive
horizontal delta. The total horizontal delta
in (C.RD.28) is one, and the two fibres are
disjoint. Away from these fibres the local
plane delta is at most one and its cubic
suspension is rational. Therefore
\[
 P:=\sum_x p_g(X_0,x)\le1,\qquad
 \lambda:=\sum_x b_1(\operatorname{link}(X_0,x))\le2.
\tag{C.RD.29}
\]

In this case the irregularity is zero, not
merely at most one. The finite-cover dualizing
formula gives
$\omega_{X_0}=p_0^*(\omega_W\otimes\OO_W(2))
\simeq\OO_{X_0}$, and (C.RD.13) gives
$H^1(X_0,\OO)=0$. Thus $\chi(\OO_{X_0})=2$.
For a resolution $\rho:X\to X_0$, local duality
and Grauert--Riemenschneider vanishing identify
$\rho_*\omega_X\subset\omega_{X_0}$ with an
ideal of colength $p_g(X_0,x)$ at each singular
point. Explicitly this is the dual of the
finite-length $R^1\rho_*\OO_X$ in the surface
duality sequence. If $P>0$, the global generator
of $\omega_{X_0}$ does not belong to that ideal
at a point with positive colength. No nonzero
global canonical form on $X$ can then exist,
because its pushforward would be a nonzero
constant multiple of that generator.
Resolution Leray now gives
\[
 \chi(\OO_X)=2-P,
 \qquad q(X)=P-1\quad(P>0),
 \qquad q(X)=0\quad(P=0).
\tag{C.RD.30}
\]
Since $P\le1$, this proves $b_1(X)=0$.

The four branch classes span the two-dimensional
space generated by a section and a fibre, so
$\varepsilon=2$ and $b_1(U)=2$.
The branch divisor and resolved boundary are
connected. For
$\nu=\sum_x(r_x(B)-1)$, normalization and the
exact pairing calculation give
\[
 e(U)=3(\nu-4),\qquad b_1(D)=\nu-3+\lambda,
 \qquad \operatorname{rank}I_U=2\nu-8-\lambda.
\tag{C.RD.31}
\]
The compact torus comparison forces this rank
to be zero; otherwise its nondegenerate compact
pairing would embed into a rank-at-most-two
first-homology image in the torus. Hence
$\nu=4+\lambda/2\le5$.

Let $a_p,a_q\in\{1,2\}$ be the numbers of
$C$-preimages of these two common values.
The exact branch count (C.RD.25) gives
$\nu_{F_p}+\nu_{F_q}=2+a_p+a_q$.
The loss of distinct attachment values on
these fibres is at most $a_p+a_q-2$.
Away from them each loss costs a multibranch
point of $R_2$. Consequently
\[
                    \ell\le\nu-4\le1.
\tag{C.RD.32}
\]
The graph count gives
$e(S)\ge2-r_{\rm dom}-\ell\ge0$.

\paragraph{One double common zero value.}
Now let $A=2p$. Normalize the cyclic cover
with branch multiplicities $R_1+R_2+2F_p$.
Its exact character algebra is
\[
 (p_0)_*\OO_{X_0}
       =\OO_W\oplus L^{-1}\oplus L^{-2}(F_p),
                  \qquad L=\OO_W(1).
\tag{C.RD.33}
\]
The additional integral generator at $F_p$
is $w^2/s$ in $w^3=s^2u$. Its divisorwise
character line extends reflexively over the
smooth base, which proves the sheaf identity
at the finitely many intersections as well;
this is also (C.DD.16). Relative duality gives
$R^1p_*L^{-2}(F_p)=\OO_H(-1)$.
Thus $H^1(X_0,\OO_{X_0})=0$.

We check that every singularity of this
normalized cover is rational, including the
horizontal node and cusp allowed by (C.RD.28).
For this purpose write a cyclic branch divisor
as $D_1+2D_2$, with root line $L$.
At a point blowup $\sigma$ let
$m_i=\operatorname{mult}D_i$,
$n=m_1+2m_2=3k+j$, $0\le j<3$, and let $E$
be the new exceptional curve. Keeping the
same function field gives root line
$L'=\sigma^*L(-kE)$ and new exceptional
branch weight $j$. Its two character sheaves
are exactly
\[
\begin{split}
 (L')^{-1}&=\sigma^*(L^{-1})(kE),\\
 (L')^{-2}(D'_2)
   &=\sigma^*(L^{-2}(D_2))
       \bigl((2k-m_2+\mathbf1_{j=2})E\bigr).
\end{split}
\tag{C.RD.34}
\]
Here $D'_2=\sigma^*D_2-m_2E+\mathbf1_{j=2}E$.
These identities follow by replacing the
cyclic root by its quotient by the $k$th
power of a local exceptional equation.

At a smooth horizontal branch the contact
with $F_p$ is one or two. The weight-$(2,1)$
contact calculation (C.DD.18)--(C.DD.20)
separates the branches with character
increments $(1,1)$ at each blowup.
If the horizontal node lies on $F_p$, there
are three smooth local branches of weights
$(2,1,1)$. They are transverse except possibly
for contact two between $F_p$ and the
quadratic branch in the case of an intersection
$R_1\cap R_2$. A blowup has $n=4,m_2=1$,
so both character increments in (C.RD.34)
are one. In the transverse case it separates
the three branches. In the contact-two case
one remaining ordinary triple has the same
weights $(2,1,1)$; the same blowup separates it.

If the horizontal cusp lies on $F_p$, its
local form is a weight-one cusp transverse
to the weight-two fibre:
\[
                        w^3=s^2(x^2-s^3).
\tag{C.RD.35}
\]
The first blowup has $n=4,m_2=1$ and
increments $(1,1)$. The strict cusp becomes
smooth with contact two with the new
weight-one exceptional curve. Blow up that
weight-$(1,1)$ tangency; here $n=2,m_2=0$
and the increments are $(0,1)$.
The remaining ordinary triple has weights
$(1,1,2)$, so a final blowup has increments
$(1,1)$. The resulting branch has normal
crossings. At a crossing of weights $a,b$
in $\{1,2\}$, normalization of
$w^3=x^a y^b$ is the quotient of $\C^2$
under the subgroup
\[
 \{(\zeta^i,\zeta^j):ai+bj=0\pmod3\},
 \qquad x=u^3,\quad y=v^3,\quad w=u^av^b.
\]
These are the rational cyclic quotient
singularities $1/3(1,1)$ or $1/3(1,2)$.
Their standard resolutions have respectively
one $(-3)$-curve or a two-curve $(-2)$ chain,
and $R^1\OO=0$. The horizontal node or cusp
away from $F_p$ likewise has a rational
cubic suspension by (C.RD.10).

For completeness, none of the preceding
base blowups adds a geometric-genus defect.
Every exceptional character increment in
(C.RD.34) was zero or one. A point blowup
has $\sigma_*\OO(E)=\OO$ and
$R^1\sigma_*\OO(E)=0$, from
$\OO(E)|_E=\OO_{\PP^1}(-1)$ and the
structure-sheaf direct images; the same
holds for increment zero. Projection formula
therefore makes the first direct image of
each normalized character algebra vanish.
The normalized cyclic covers over successive
base blowups map properly and birationally
to the preceding cover through their same
function field. Finite direct image and
Leray identify their $R^1\OO$ with the
just computed eigensheaf direct images;
faithfulness of finite pushforward gives
$R^1\OO=0$. Resolving the final rational
cyclic quotients preserves this property.
This proves the claimed local rationality.

We consequently have $b_1(X)=0$ and
$\lambda=0$. The reduced branch support
$B=R_1+R_2+F_p$ has three rational
normalization components, is connected,
and has class rank two. Thus
\[
 b_1(U)=1,\qquad e(U)=3(\nu-2),\qquad
 b_1(D)=\nu-2,\qquad
          \operatorname{rank}I_U=2\nu-4.
\tag{C.RD.36}
\]
Again compact transport forces rank zero,
hence $\nu=2$.
Equation (C.RD.25) then gives exactly one
$C$-preimage over the common value and no
multibranch point of $R_2$ away from it.
There is no loss of distinct attachment
values. The graph inequality yields
\[
                    e(S)\ge2-r_{\rm dom}\ge1.
\tag{C.RD.37}
\]
This proves the degree-two-common-divisor
case and, together with C.2.RAT.7--8 and
the degree bound at the start, every case
in the statement.
The additional weighted node/cusp charts,
their exact character increments and both
attachment budgets are checked in
\path{code/A_cubic_scalar_two.py} and its output.
\end{proof}


% C.2.UN.61 from part_C2_rooted_hopf_trees.tex

